% ===================================================================
\section{Structured populations: life tables, matrices, elasticity (EE02)}
% ===================================================================

\subsection{What a population is and what we ask of it}
\textbf{Ecology} (Haeckel, 1866; \emph{oikos} = house) is the study of interactions among organisms and their biotic
and abiotic environment, and how these lead to patterns of distribution and abundance. A \textbf{population} is a
group of similar individuals inhabiting a defined space at a given time. Population ecologists study
\textbf{distribution} (spread over space), \textbf{demography} (age structure, survivorship, birth rates) and
\textbf{dynamics} (how and why numbers change). Reasons to study non-human populations: conservation, management of
useful species, control of pests and invaders, and the population being the \textbf{unit of evolution}.

The population is a ledger: $N_{t+1}=N_t+B-D+I-E$. If every individual had the same per-capita rates, one number
would describe growth. This chapter asks what happens when \textbf{individuals are not interchangeable}: a newborn, a
breeder and an old animal survive and reproduce differently. That is \emph{population structure}.

\subsection{The cohort life table and $R_0$}
Follow a \textbf{cohort} (individuals born in the same period) to the last death. Columns: $n_x$ alive at the start of
class $x$; $l_x=n_x/n_0$ the proportion surviving to $x$; $m_x$ the mean offspring per survivor in class $x$; and
$l_xm_x$ the expected offspring per \emph{newborn} from that class. Summing gives the \textbf{net reproductive rate}
\[
R_0=\sum_x l_x m_x ,
\]
the mean number of offspring a newborn produces over its whole life. $R_0>1$: each generation larger than the last;
$R_0<1$: shrinking. \emph{Phlox drummondii} (annual plant, cohort of 996): $R_0=2.4177$. Its survival is front-loaded
(84\,\% die before day 299) and reproduction concentrated in a few late weeks.

Three ways to get survival data (Molles): \textbf{cohort} table (most reliable, impractical for long-lived or mobile
species); \textbf{static} table (ages at death of many individuals, e.g.\ Dall sheep skulls; assumes the death-age
distribution represents survival); \textbf{from an age distribution} (differences between classes; assumes a
stationary, closed population, so least accurate).

\begin{trap}
$R_0$ is a factor \emph{per generation}; $\lambda$ is a factor \emph{per time step}. They are equal only for
non-overlapping generations (an annual plant). In general they agree only on the threshold: $\lambda>1\iff R_0>1$.
\end{trap}

Unlimited growth at $\lambda=2.42$ would cover North America with $\sim$77 \emph{Phlox} per m$^2$ in 32 years:
natural populations have a huge capacity for increase, so something must change as $N$ grows (density dependence,
next section).

\subsection{Survivorship curves and the log axis}
\begin{chain}[: why survivorship is plotted on a log axis]
If an individual has the same probability $q$ of dying in every interval, then $n_{x+1}=(1-q)n_x$, so
$n_x=n_0(1-q)^x$ and
\[
\ln n_x=\ln n_0+x\ln(1-q).
\]
Constant per-capita risk is therefore a \textbf{straight line} on a log axis, with slope $\ln(1-q)$. Slide 12: 60\,\%
die each interval; numbers 1000, 400, 160, 64. The \emph{losses} shrink (600, 240, 96) but the \emph{proportion}
lost (and the log-slope) is the same each time. A log axis shows the risk to \emph{one individual}; a linear axis
shows the number of bodies. Population dynamics asks about per-capita rates, so the log axis is the right one.
\end{chain}

\begin{itemize}
\item \textbf{Type I}: low mortality until late in life (large mammals, humans; Dall sheep, \emph{Phlox} after its
first period).
\item \textbf{Type II}: constant proportion dies per interval; straight line on a log axis (many birds: robin,
white-crowned sparrow).
\item \textbf{Type III}: very high juvenile mortality, then high survival (marine fish, mackerel; the desert shrub
\emph{Cleome}).
\end{itemize}
The idealised types are \emph{bounds} on what is possible; real curves are mixtures (\emph{Phlox}: a Type III-like
early loss, a plateau, then a Type I-like end).

\begin{trap}
A steep early drop (Type III) says nothing about lifespan: the survivors may live long. And the slope on a log axis is
the per-capita hazard, not the number of deaths.
\end{trap}

\subsection{The Leslie matrix: age structure with births}
A cohort table ignores the cohort's children. To project the whole population, count every age class and let births
join class 1:
\[
\begin{pmatrix}n_1\\n_2\\n_3\\n_4\end{pmatrix}_{t+1}=
\begin{pmatrix}F_1&F_2&F_3&F_4\\P_1&0&0&0\\0&P_2&0&0\\0&0&P_3&0\end{pmatrix}
\begin{pmatrix}n_1\\n_2\\n_3\\n_4\end{pmatrix}_{t}.
\]
Row 1 reads ``babies = $\sum$ fecundity $\times$ number in each class''; each sub-diagonal row reads ``class $x+1$ =
survivors of class $x$''.

Slide 16: $P=(0.5,0.2,0.1)$, $F=(0,2,1,1)$, start with 200 newborns. The totals first zig-zag
(200, 100, 220, 122, 232, \dots) because one cohort breeds in pulses; then
\begin{itemize}
\item the total grows \textbf{geometrically} by $\lambda=1.05082$ per step (slide 18, 20);
\item the age \emph{proportions} settle at $(0.635,\,0.302,\,0.057,\,0.005)$: the \textbf{stable age distribution}
(SAD), even though the numbers keep growing (slide 19).
\end{itemize}

\begin{chain}[: why the population settles (eigenvalues)]
\begin{enumerate}
\item Try a solution in which every class grows by the same factor: $\mathbf n_t=\lambda^t\mathbf w$. Then
$\mathbf A\mathbf w=\lambda\mathbf w$: $\lambda$ is an eigenvalue and $\mathbf w$ its eigenvector.
\item Any start is a sum of such modes, $\sum_i c_i\lambda_i^t\mathbf w_i$. After many steps the mode with the
largest $|\lambda_i|$ dominates; the others shrink \emph{relative to it} like $(|\lambda_i|/\lambda_1)^t$.
\item So the shape converges to $\mathbf w_1$ (the SAD) and the growth ratio to $\lambda_1$, \emph{whatever the
starting mix}. The matrix ``forgets'' its start (ergodicity). The SAD depends only on the vital rates, not the initial
vector.
\item The sub-diagonal rows give $w_{x+1}/w_x=P_x/\lambda$, and the first row gives the discrete Euler--Lotka equation
\[
1=\sum_x F_x\,l_x\,\lambda^{-x}.
\]
Its right side falls steadily as $\lambda$ rises, so it has one positive root; at $\lambda=1$ it equals $R_0$. Hence
$\lambda>1\iff R_0>1$.
\end{enumerate}
\end{chain}

\textbf{Perron--Frobenius.} The slide states that a non-negative matrix has a real dominant eigenvalue. For
convergence to a fixed shape the matrix must be non-negative, \textbf{irreducible} (life-cycle graph connected) and
\textbf{primitive} (the loop lengths have greatest common divisor 1). Then $\lambda_1$ is real, positive, simple and
strictly larger in modulus than every other eigenvalue, and $\mathbf w$ (SAD) and $\mathbf v$ (reproductive values)
are strictly positive.

\begin{trap}
A species that reproduces in only \emph{one} age class (semelparous, single loop) gives an \emph{imprimitive} matrix:
e.g.\ $F_3=6$, $P_1=0.5$, $P_2=0.4$ repeats a 3-step cycle (100, 50, 20, 120, 60, 24, \dots) growing 1.2$\times$
per cycle, with no SAD and no limit for $N_{t+1}/N_t$. ``Every Leslie matrix converges to a SAD'' is false without
primitivity.
\end{trap}

\subsection{Stage structure: the loggerhead turtle}
For turtles, \textbf{age is hard to measure} and \textbf{stage (size) decides fate}. So classify by stage and let an
individual stay ($P$), grow to the next stage ($G$) or die (Lefkovitch matrix). The pedagogical move (slide 30):
start oversimplified (cohort table), remove one constraint at a time (births, then stages).

\begin{center}
\small
$\mathbf A=\begin{pmatrix}0&0&0&4.665&61.896\\.675&.703&0&0&0\\0&.047&.657&0&0\\0&0&.019&.682&0\\0&0&0&.061&.809\end{pmatrix}$
\qquad
\begin{tabular}{lcc}
\toprule
stage & years & survival $\sigma$\\\midrule
1 eggs/hatchlings & 1 & 0.675\\
2 small juveniles & 7 & 0.75\\
3 large juveniles & 8 & 0.676\\
4 subadults & 6 & 0.743\\
5 adults & $>$32 & 0.809 (76.5 eggs/yr)\\\bottomrule
\end{tabular}
\end{center}

$P$ and $G$ each include survival: $P_i+G_i=\sigma_i$ (e.g.\ $0.703+0.047=0.750$); mortality is what is left out.
A long stage has a \emph{small} $G$ because only its oldest members graduate each year
($G=\sigma\gamma$, $\gamma=x^{T-1}(1-x)/(1-x^T)$, $x=\sigma/\lambda$).

\begin{itemize}
\item $\lambda=0.9516$: decline of about \textbf{5\,\% a year}.
\item Stable stage distribution: small juveniles $\sim$65\,\%, adults only 0.3\,\%.
\item Reproductive values (hatchling $=1$): 1, 1.4, 7.5, 115, 432. The jump is largest across the
\textbf{large-juvenile} stage.
\end{itemize}

\subsubsection*{Where should conservation act? Sensitivity and elasticity}
\[
s_{ij}=\frac{\partial\lambda}{\partial a_{ij}}=\frac{v_iw_j}{\mathbf v\cdot\mathbf w},\qquad
e_{ij}=\frac{a_{ij}}{\lambda}\,s_{ij},\qquad \sum_{ij}e_{ij}=1.
\]
Sensitivity is the \emph{absolute} change of $\lambda$ per absolute change of an entry; elasticity the
\emph{proportional} change per proportional change. The formula says $\lambda$ responds most to a transition that
moves many individuals ($w_j$ large) into a valuable state ($v_i$ large).

\begin{chain}[: why egg protection cannot save the loggerhead]
\begin{enumerate}
\item The egg/hatchling stage lasts one year and leads to small juveniles worth only 1.4 hatchlings: low $v$,
one chance to act.
\item Large juveniles stay 8 years (eight chances per animal to be lost) and each survivor is about to jump $\sim$15-fold
in reproductive value. Their survival has the highest elasticity.
\item Recomputed (slide 26, Crowder et al.\ Fig.~4): cutting mortality by 50\,\% makes $r>0$ only for large juveniles;
even \emph{perfect} egg survival leaves $\lambda=0.965<1$, while perfect large-juvenile survival gives $\lambda=1.073$.
\item What kills large juveniles? \textbf{Shrimp trawling} (5.7\,kg by-catch per kg shrimp; large juveniles are the
commonest stranded size). So the intervention is \textbf{turtle excluder devices} (TEDs), escape hatches in nets.
\item Recovery is slow because maturity takes $>$20 years: with seasonal-offshore TED use, $\lambda\approx1.035$ and a
ten-fold increase takes $\sim$70 years; with full use, about half that. Hence the advice to use TEDs all the time
(Indian legislation 1996, per the slide).
\end{enumerate}
\end{chain}

\begin{trap}
\emph{Largest sensitivity} is not the same as \emph{most useful target}. $G_3$ has the largest sensitivity (2.9)
because the entry is tiny (0.019); its elasticity is small. Use elasticity to rank, then ask what is feasible.
Also: a flat nest count a decade after TEDs is not evidence they failed (background: transient ``waves'' as the
age structure adjusts).
\end{trap}

\subsection{Age pyramids}
Slide 29: \textbf{the greater the $\lambda$, the broader the base}. Reason: at the SAD $w_{x+1}/w_x=P_x/\lambda$.
Each older class is smaller because (i) some of it died ($P_x<1$) and (ii) it was born when the population was smaller
(the factor $1/\lambda$). In a growing population the second factor narrows each older class further, giving a broad
base; in a declining one, older classes were born into a larger population, so the base is narrow.
Survival sets the slope of the sides; $\lambda$ tilts them.

But real populations are rarely at their SAD, so a pyramid also records \textbf{history}: white oak (many young trees:
replacing itself), Rio Grande cottonwood (dominated by 40--50-year-olds, no recruitment because dams stopped the
floods seedlings need), the large cactus finch on Genovesa (gaps from drought years).

\begin{numbers}
Sweden $r=0$; Hungary $r=-0.003$ (halving 231\,yr); Rwanda $r=+0.027$ (doubling $\ln2/r=25.7$\,yr).
Slide-16 Leslie example: $\lambda=1.05082$, SAD (0.635, 0.302, 0.057, 0.005). Loggerhead $\lambda=0.9516$;
by-catch 5.7\,kg/kg shrimp; TEDs: seasonal $\approx$70\,yr for 10$\times$, full use $\approx$33\,yr.
\end{numbers}

\begin{statements}
\begin{enumerate}[label=\arabic*.]
\item ``A broad-based pyramid always means a growing population.''
\PT At the SAD, yes. Away from it, a broad base can reflect very high juvenile mortality (many born, few reach
breeding age), a one-off pulse of young (a good year, juvenile immigration), or removal of adults (harvest, disease)
that cuts the top; the pyramid records history as well as $\lambda$ (2022 mid-sem Q4).
\item ``The stable age distribution depends on the initial age structure.''
\F It is the dominant eigenvector of the matrix; only the transient depends on the start.
\item ``At the SAD the number in each class is constant.''
\F The \emph{proportions} are constant; every class grows by $\lambda$ per step.
\item ``A type II curve means the same number of individuals die each year.''
\F The same \emph{proportion} dies; the numbers dying shrink.
\item ``Since adults have the highest reproductive value, protecting adults is the only way to stop the loggerhead decline.''
\F Large-juvenile survival has the higher elasticity and is where trawls kill; adults are only 0.3\,\% of the
population.
\item ``If $R_0>1$ then $\lambda>1$, though the two numbers generally differ.''
\T Euler--Lotka gives $\lambda>1\iff R_0>1$; they coincide only with non-overlapping generations.
\item ``An analyst sees an urn-shaped pyramid (narrow base, bulge of middle-aged adults) in a low-birth country and predicts a demographic dividend in 20 years.''
\F In 20 years the bulge will be old and dependent, and the narrow base means fewer new workers: a falling worker
ratio (2025 end-sem Q2).
\item ``A slot limit (only medium-sized fish may be caught) is more sustainable than a minimum-size limit.''
\T It spares the large, old fish with the highest fecundity, the stage to which $\lambda$ is most sensitive in
fish, and avoids selecting for early maturation at small size (2025 end-sem Q4).
\end{enumerate}
\end{statements}

\begin{longans}[: ``On which stage should loggerhead conservation be focused, and why?'']
(1) Set up the stage matrix ($P$, $G$, $F$; $P+G=\sigma$) and state $\lambda=0.95$. (2) Define sensitivity and
elasticity; explain why elasticity ranks entries of different sizes fairly. (3) Large juveniles: long stage, high
reproductive value ahead, highest elasticity; show eggs alone cannot give $\lambda>1$. (4) Cause of their mortality:
shrimp trawls $\Rightarrow$ TEDs. (5) Time scale: $>$20 years to mature, so 33--70 years for ten-fold recovery;
counts may dip in between. (6) Caveat: the model is deterministic, data poor, and assumes TED savings are additive;
use it to rank parameters, not to forecast precisely.
\end{longans}
